\documentclass[11pt,a4paper]{article}

\usepackage[margin=2.2cm]{geometry}
\usepackage{kotex}
\usepackage{amsmath,amssymb}
\usepackage{booktabs}
\usepackage{array}
\usepackage{xcolor}
\usepackage{listings}
\usepackage[most]{tcolorbox}
\usepackage{enumitem}
\usepackage{hyperref}
\usepackage{fancyhdr}
\usepackage{titlesec}

\definecolor{codebg}{RGB}{247,247,247}
\definecolor{codeframe}{RGB}{210,210,210}
\definecolor{keywordcolor}{RGB}{0,70,160}
\definecolor{commentcolor}{RGB}{0,120,80}
\definecolor{stringcolor}{RGB}{170,50,50}
\definecolor{titleblue}{RGB}{35,75,130}

\hypersetup{colorlinks=true,linkcolor=titleblue,urlcolor=titleblue}

\pagestyle{fancy}
\fancyhf{}
\lhead{C++ Programming}
\rhead{Day 9}
\cfoot{\thepage}

\titleformat{\section}{\Large\bfseries\color{titleblue}}{\thesection.}{0.6em}{}
\titleformat{\subsection}{\large\bfseries}{\thesubsection.}{0.5em}{}

\lstdefinestyle{cppstyle}{
    language=C++,
    backgroundcolor=\color{codebg},
    frame=single,
    rulecolor=\color{codeframe},
    basicstyle=\ttfamily\small,
    keywordstyle=\color{keywordcolor}\bfseries,
    commentstyle=\color{commentcolor},
    stringstyle=\color{stringcolor},
    numbers=left,
    numberstyle=\tiny\color{gray},
    numbersep=8pt,
    tabsize=4,
    showstringspaces=false,
    breaklines=true,
    columns=fullflexible
}
\lstset{style=cppstyle}

\newtcolorbox{learningbox}{colback=blue!4,colframe=titleblue,title=\textbf{학습 목표},fonttitle=\bfseries}
\newtcolorbox{importantbox}{colback=yellow!8,colframe=orange!70!black,title=\textbf{핵심 포인트},fonttitle=\bfseries}
\newtcolorbox{practicebox}{colback=red!3,colframe=red!55!black,title=\textbf{실습},fonttitle=\bfseries}

\begin{document}

\begin{titlepage}
    \centering
    \vspace*{3cm}
    {\Huge\bfseries C++ Programming\par}
    \vspace{0.8cm}
    {\LARGE Day 9: 템플릿\par}
    \vspace{2cm}
    {\large 함수 템플릿, 다중 타입, 오버로딩, 클래스 템플릿, 기본 템플릿 인자\par}
    \vfill
\end{titlepage}

\tableofcontents
\newpage

\section{학습 목표}
\begin{learningbox}
이 장을 마치면 다음 내용을 설명하고 직접 코드로 작성할 수 있다.
\begin{itemize}[leftmargin=1.5em]
    \item 코드 중복이 발생하는 이유와 템플릿의 필요성 설명하기
    \item \texttt{template <typename T>} 문법으로 함수 템플릿 작성하기
    \item 호출 인수로부터 템플릿 자료형이 추론되는 과정 설명하기
    \item 여러 개의 템플릿 타입 매개변수를 동시에 사용하기
    \item 일반 함수와 템플릿 함수가 함께 있을 때 호출 우선순위 이해하기
    \item \texttt{template <typename T>} 문법으로 클래스 템플릿 작성하기
    \item 객체를 만들 때 클래스 템플릿의 실제 자료형 지정하기
    \item 기본 템플릿 인자를 선언하고 생략 문법 사용하기
    \item 템플릿 자료형 결정과 암시적 형 변환을 구분하기
    \item 템플릿이 이후 STL 컨테이너의 기반이 되는 이유 설명하기
\end{itemize}
\end{learningbox}

\section{템플릿의 필요성}
자료형만 다르고 동작은 같은 함수나 클래스를 여러 번 작성하면 코드 중복이 발생한다. 다음 두 함수는 매개변수와 반환형만 다르고 수행하는 연산은 같다.

\begin{lstlisting}
int add(int a, int b)
{
    return a + b;
}

double add(double a, double b)
{
    return a + b;
}
\end{lstlisting}

자료형이 \texttt{float}, \texttt{long}, \texttt{string}으로 늘어나면 같은 형태의 코드를 계속 반복해야 한다. 템플릿(template)은 자료형을 미리 고정하지 않고, 실제 호출이나 객체 생성 시점에 필요한 자료형을 컴파일러가 결정하도록 만드는 기능이다.

\begin{verbatim}
Repeated code
int add(int, int)
double add(double, double)
float add(float, float)

        |
        v

One template
T add(T, T)
\end{verbatim}

템플릿은 실행 중에 자료형이 바뀌는 기능이 아니다. 컴파일러가 코드를 번역할 때 실제로 사용된 자료형에 맞는 함수 또는 클래스를 생성한다.

\begin{importantbox}
템플릿은 자료형에 의존하지 않는 공통 알고리즘을 한 번만 작성하게 한다. 실제로 사용되는 자료형별 코드는 컴파일 시점에 생성된다.
\end{importantbox}

\section{함수 템플릿}
함수 템플릿(function template)은 자료형만 다른 여러 함수를 하나의 일반화된 형태로 작성하는 기능이다.

\subsection{기본 문법}
함수 선언 앞에 다음 문장을 작성한다.

\begin{lstlisting}
template <typename T>
\end{lstlisting}

\texttt{T}는 아직 결정되지 않은 자료형을 나타내는 템플릿 타입 매개변수(template type parameter)이다. 이름은 반드시 \texttt{T}일 필요는 없지만, 관례적으로 첫 번째 타입에 \texttt{T}를 자주 사용한다.

\begin{lstlisting}
template <typename T>
T add(T a, T b)
{
    return a + b;
}
\end{lstlisting}

위 코드에서 \texttt{T}는 반환형, 첫 번째 매개변수형, 두 번째 매개변수형에 모두 사용된다.

\subsection{typename과 class}
템플릿 타입 매개변수 선언에는 다음 두 문법을 사용할 수 있다.

\begin{lstlisting}
template <typename T>
\end{lstlisting}

\begin{lstlisting}
template <class T>
\end{lstlisting}

기본적인 타입 매개변수 선언에서는 두 표현의 의미가 같다. 현대 C++ 코드에서는 해당 매개변수가 자료형임을 더 직접적으로 나타내는 \texttt{typename}을 자주 사용한다.

\subsection{자료형 추론}
함수 템플릿을 호출하면 컴파일러가 전달된 인수의 자료형을 보고 \texttt{T}를 추론한다.

\begin{lstlisting}
add(10, 20);
\end{lstlisting}

두 인수가 모두 \texttt{int}이므로 컴파일러는 개념적으로 다음 함수를 만든다.

\begin{lstlisting}
int add(int a, int b)
{
    return a + b;
}
\end{lstlisting}

다음 호출에서는 \texttt{T}가 \texttt{double}이 된다.

\begin{lstlisting}
add(2.5, 3.1);
\end{lstlisting}

\begin{lstlisting}
double add(double a, double b)
{
    return a + b;
}
\end{lstlisting}

이처럼 하나의 템플릿 정의로 여러 자료형에 대응하는 함수를 생성할 수 있다.

\subsection{같은 T가 의미하는 것}
다음 템플릿에서 두 매개변수는 모두 같은 \texttt{T}를 사용한다.

\begin{lstlisting}
template <typename T>
T multiply(T a, T b)
\end{lstlisting}

따라서 다음 호출은 자연스럽게 추론된다.

\begin{lstlisting}
multiply(3, 4);
multiply(2.5, 4.2);
\end{lstlisting}

그러나 다음 호출에서는 첫 번째 인수가 \texttt{int}, 두 번째 인수가 \texttt{double}이므로 하나의 \texttt{T}를 결정할 수 없다.

\begin{lstlisting}
multiply(3, 2.5);
\end{lstlisting}

이 문제는 서로 다른 타입 매개변수를 두 개 선언하면 해결할 수 있다.

\subsection{전체 예제}
\begin{lstlisting}
#include <iostream>

using namespace std;

template <typename T>
T add(T a, T b)
{
    return a + b;
}

int main()
{
    cout << add(10, 20) << endl;
    cout << add(2.5, 3.1) << endl;

    return 0;
}
\end{lstlisting}

출력 결과는 다음과 같다.

\begin{verbatim}
30
5.6
\end{verbatim}

\section{여러 개의 템플릿 타입 매개변수}
서로 다른 자료형을 동시에 받아야 한다면 템플릿 타입 매개변수를 여러 개 선언한다.

\begin{lstlisting}
template <typename T, typename U>
\end{lstlisting}

\texttt{T}와 \texttt{U}는 서로 독립적인 자료형이다. 첫 번째 인수에서 \texttt{T}를 추론하고 두 번째 인수에서 \texttt{U}를 추론할 수 있다.

\subsection{독립적인 자료형 추론}
다음 호출을 살펴보자.

\begin{lstlisting}
show(100, 3.14);
\end{lstlisting}

컴파일러는 다음과 같이 추론한다.

\begin{verbatim}
T = int
U = double
\end{verbatim}

다음 호출에서는 다른 조합이 만들어진다.

\begin{lstlisting}
show("Age", 20);
\end{lstlisting}

문자열 리터럴은 함수 매개변수로 전달될 때 일반적으로 문자열을 가리키는 포인터 형태로 처리되므로 개념적으로 다음과 같은 조합이 된다.

\begin{verbatim}
T = const char*
U = int
\end{verbatim}

\subsection{전체 예제}
\begin{lstlisting}
#include <iostream>

using namespace std;

template <typename T, typename U>
void show(T first, U second)
{
    cout << first << " " << second << endl;
}

int main()
{
    show(100, 3.14);
    show("Age", 20);

    return 0;
}
\end{lstlisting}

출력 결과는 다음과 같다.

\begin{verbatim}
100 3.14
Age 20
\end{verbatim}

\subsection{타입 매개변수의 개수}
필요하다면 세 개 이상의 타입 매개변수도 선언할 수 있다.

\begin{lstlisting}
template <typename T, typename U, typename V>
void printThree(T first, U second, V third)
{
    cout << first << " " << second << " " << third << endl;
}
\end{lstlisting}

다만 타입 매개변수가 많아질수록 함수의 사용 조건과 반환형 설계가 복잡해질 수 있으므로 실제 필요에 맞게 사용한다.

\begin{importantbox}
같은 이름의 타입 매개변수는 같은 자료형을 의미한다. 서로 다른 자료형을 독립적으로 받으려면 \texttt{T}, \texttt{U}처럼 별도의 타입 매개변수를 선언해야 한다.
\end{importantbox}

\section{일반 함수와 템플릿 함수}
일반 함수와 함수 템플릿은 같은 이름으로 함께 존재할 수 있다. 이때 컴파일러는 함수 오버로딩 규칙에 따라 가장 적합한 후보를 선택한다.

\subsection{정확한 일치와 형 변환}
다음 두 함수가 있다고 하자.

\begin{lstlisting}
void print(int x)
{
    cout << "Normal Function" << endl;
}

template <typename T>
void print(T x)
{
    cout << "Template Function" << endl;
}
\end{lstlisting}

\texttt{print(10)}은 일반 함수 \texttt{print(int)}와 정확히 일치한다. 템플릿으로도 \texttt{print(int)}를 만들 수 있지만, 동일하게 정확한 후보라면 일반 함수가 선택된다.

\begin{lstlisting}
print(10);
\end{lstlisting}

\begin{verbatim}
Normal Function
\end{verbatim}

반면 \texttt{print(3.14)}에서 일반 함수는 \texttt{double}을 \texttt{int}로 바꾸는 형 변환이 필요하다. 템플릿은 \texttt{T = double}로 정확히 일치하는 함수를 만들 수 있으므로 템플릿 함수가 선택된다.

\begin{lstlisting}
print(3.14);
\end{lstlisting}

\begin{verbatim}
Template Function
\end{verbatim}

\subsection{선택 원리}
함수 선택은 단순히 ``일반 함수가 항상 우선''이라고 이해하면 안 된다. 먼저 각 후보가 인수와 얼마나 잘 일치하는지 비교한다.

\begin{center}
\begin{tabular}{ll}
\toprule
상황 & 선택 결과 \\
\midrule
일반 함수만 정확히 일치 & 일반 함수 \\
템플릿 함수만 정확히 일치 & 템플릿 함수 \\
둘 다 정확히 일치 & 일반 함수 \\
일반 함수는 형 변환 필요, 템플릿은 정확히 일치 & 템플릿 함수 \\
\bottomrule
\end{tabular}
\end{center}

\subsection{전체 예제}
\begin{lstlisting}
#include <iostream>

using namespace std;

void print(int x)
{
    cout << "Normal Function" << endl;
}

template <typename T>
void print(T x)
{
    cout << "Template Function" << endl;
}

int main()
{
    print(10);
    print(3.14);

    return 0;
}
\end{lstlisting}

출력 결과는 다음과 같다.

\begin{verbatim}
Normal Function
Template Function
\end{verbatim}

\subsection{double 일반 함수와의 비교}
다음 코드에서는 첫 번째 호출과 두 번째 호출의 선택 결과가 다르다.

\begin{lstlisting}
void print(double x)
{
    cout << "Double Function" << endl;
}

template <typename T>
void print(T x)
{
    cout << "Template Function" << endl;
}
\end{lstlisting}

\texttt{print(3.14)}은 일반 함수가 정확히 일치하므로 \texttt{Double Function}이 출력된다. \texttt{print(100)}은 일반 함수에서 \texttt{int}를 \texttt{double}로 바꿔야 하지만, 템플릿은 \texttt{T = int}로 정확히 일치하므로 \texttt{Template Function}이 출력된다.

\begin{importantbox}
정확한 일치 여부가 먼저 비교된다. 일반 함수와 템플릿 함수가 모두 같은 수준으로 정확히 일치할 때 일반 함수가 선택된다.
\end{importantbox}

\section{클래스 템플릿}
클래스 템플릿(class template)은 멤버 변수나 멤버 함수가 사용할 자료형을 객체 생성 시점까지 미루는 기능이다.

\subsection{클래스 중복 문제}
값 하나를 저장하는 클래스를 자료형별로 만들면 다음과 같이 코드가 반복된다.

\begin{lstlisting}
class IntBox
{
public:
    int value;
};

class DoubleBox
{
public:
    double value;
};
\end{lstlisting}

클래스 템플릿을 사용하면 하나의 클래스 정의로 여러 자료형을 처리할 수 있다.

\begin{lstlisting}
template <typename T>
class Box
{
public:
    T value;
};
\end{lstlisting}

\subsection{객체 생성 시 자료형 지정}
함수 템플릿은 인수로부터 타입을 추론할 수 있지만, 기본적인 클래스 템플릿 사용에서는 객체를 만들 때 각괄호 안에 실제 자료형을 명시한다.

\begin{lstlisting}
Box<int> box1;
Box<double> box2;
Box<char> box3;
\end{lstlisting}

\texttt{Box<int>}와 \texttt{Box<double>}은 같은 템플릿에서 만들어지지만 서로 다른 구체 클래스 타입이다.

\begin{verbatim}
Box<T>
  |
  +-- Box<int>
  +-- Box<double>
  +-- Box<char>
\end{verbatim}

\subsection{생성자와 멤버 함수}
클래스 내부에서는 템플릿 타입 매개변수 \texttt{T}를 멤버 변수, 생성자 매개변수, 반환형 등 필요한 위치에 사용할 수 있다.

\begin{lstlisting}
template <typename T>
class Box
{
public:
    T value;

    Box(T value)
    {
        this->value = value;
    }

    void show()
    {
        cout << value << endl;
    }
};
\end{lstlisting}

\texttt{Box<int>}를 만들면 \texttt{T value}와 \texttt{Box(T value)}의 \texttt{T}가 모두 \texttt{int}로 바뀐다. \texttt{Box<double>}에서는 모두 \texttt{double}로 바뀐다.

\subsection{전체 예제}
\begin{lstlisting}
#include <iostream>

using namespace std;

template <typename T>
class Box
{
public:
    T value;

    Box(T value)
    {
        this->value = value;
    }

    void show()
    {
        cout << value << endl;
    }
};

int main()
{
    Box<int> box1(100);
    Box<double> box2(3.14);
    Box<char> box3('A');

    box1.show();
    box2.show();
    box3.show();

    return 0;
}
\end{lstlisting}

출력 결과는 다음과 같다.

\begin{verbatim}
100
3.14
A
\end{verbatim}

\subsection{STL과의 연결}
표준 템플릿 라이브러리(STL)의 컨테이너는 클래스 템플릿으로 구현되어 있다. 다음 문법의 각괄호 안 자료형은 클래스 템플릿의 타입 인자이다.

\begin{lstlisting}
vector<int> numbers;
vector<double> values;
stack<char> characters;
map<string, int> scores;
\end{lstlisting}

따라서 클래스 템플릿을 이해하면 이후 \texttt{vector}, \texttt{list}, \texttt{map}, \texttt{stack}, \texttt{queue} 등의 문법을 자연스럽게 해석할 수 있다.

\begin{importantbox}
클래스 템플릿은 하나의 클래스가 여러 자료형을 저장하거나 처리하도록 만든다. \texttt{Box<int>}와 \texttt{Box<double>}은 같은 템플릿에서 생성되지만 서로 다른 타입이다.
\end{importantbox}

\section{기본 템플릿 인자}
기본 템플릿 인자(default template argument)는 사용자가 자료형을 명시하지 않았을 때 사용할 기본 자료형을 지정하는 기능이다.

\subsection{기본 문법}
다음 템플릿은 \texttt{T}의 기본 자료형을 \texttt{int}로 지정한다.

\begin{lstlisting}
template <typename T = int>
class Number
{
};
\end{lstlisting}

객체를 만들 때 빈 각괄호를 사용하면 기본 자료형이 적용된다.

\begin{lstlisting}
Number<> n1;
\end{lstlisting}

위 코드는 다음 코드와 같은 의미이다.

\begin{lstlisting}
Number<int> n1;
\end{lstlisting}

직접 자료형을 명시하면 기본값은 사용되지 않는다.

\begin{lstlisting}
Number<double> n2;
\end{lstlisting}

\subsection{전체 예제}
\begin{lstlisting}
#include <iostream>

using namespace std;

template <typename T = int>
class Number
{
public:
    T value;

    Number(T value)
    {
        this->value = value;
    }

    void show()
    {
        cout << value << endl;
    }
};

int main()
{
    Number<> n1(100);
    Number<double> n2(3.14);

    n1.show();
    n2.show();

    return 0;
}
\end{lstlisting}

출력 결과는 다음과 같다.

\begin{verbatim}
100
3.14
\end{verbatim}

\subsection{템플릿 인자 결정과 형 변환}
다음 클래스에서 기본 자료형은 \texttt{double}이다.

\begin{lstlisting}
template <typename T = double>
class Value
{
public:
    T value;

    Value(T value)
    {
        this->value = value;
    }
};
\end{lstlisting}

다음 객체 생성에서 전달한 \texttt{10}은 \texttt{int} 리터럴이다.

\begin{lstlisting}
Value<> v1(10);
\end{lstlisting}

그러나 처리 순서는 다음과 같다.

\begin{verbatim}
1. Value<> uses the default template argument.
2. T becomes double.
3. The constructor becomes Value(double value).
4. The int value 10 is converted to double.
5. The member stores 10.0.
\end{verbatim}

즉, \texttt{10}의 자료형을 보고 \texttt{T}를 다시 \texttt{int}로 정하는 것이 아니다. \texttt{Value<>}에서 기본 템플릿 인자가 먼저 적용되어 \texttt{Value<double>}이 되고, 그 다음 생성자 호출 과정에서 \texttt{int}가 \texttt{double}로 암시적 형 변환된다.

반대 방향도 문법적으로 가능하다.

\begin{lstlisting}
Value<int> v2(3.14);
\end{lstlisting}

이 경우 \texttt{T}는 명시적으로 \texttt{int}이며, \texttt{3.14}가 \texttt{int}로 변환되어 소수 부분이 제거된 \texttt{3}이 저장된다.

\begin{center}
\begin{tabular}{lll}
\toprule
코드 & 결정된 T & 생성자 전달 과정 \\
\midrule
\texttt{Value<> v1(10)} & \texttt{double} & \texttt{int}에서 \texttt{double}로 변환 \\
\texttt{Value<char> v2('Z')} & \texttt{char} & 정확히 일치 \\
\texttt{Value<int> v3(3.14)} & \texttt{int} & \texttt{double}에서 \texttt{int}로 변환 \\
\bottomrule
\end{tabular}
\end{center}

\begin{importantbox}
템플릿 자료형을 결정하는 단계와 생성자 인수의 형 변환 단계는 서로 다르다. \texttt{Value<>}에서는 기본 템플릿 인자가 먼저 적용되고, 그 뒤 전달값이 생성자의 매개변수형에 맞게 변환된다.
\end{importantbox}

\section{종합 정리}
\begin{importantbox}
\begin{itemize}[leftmargin=1.5em]
    \item 템플릿은 자료형만 다른 중복 코드를 하나의 일반화된 코드로 작성하게 한다.
    \item 템플릿은 런타임 기능이 아니라 컴파일 시점에 필요한 코드를 생성하는 기능이다.
    \item 함수 템플릿은 \texttt{template <typename T>} 뒤에 함수를 정의한다.
    \item 함수 호출 인수로부터 템플릿 타입을 자동 추론할 수 있다.
    \item 같은 \texttt{T}를 사용한 여러 매개변수는 같은 자료형으로 추론되어야 한다.
    \item 서로 다른 자료형을 받으려면 \texttt{T}, \texttt{U}처럼 여러 타입 매개변수를 선언한다.
    \item 일반 함수와 템플릿 함수의 선택에서는 정확한 일치 여부가 중요하다.
    \item 일반 함수와 템플릿 함수가 모두 정확히 일치하면 일반 함수가 선택된다.
    \item 클래스 템플릿은 객체 생성 시 실제 자료형을 각괄호 안에 지정한다.
    \item 기본 템플릿 인자는 자료형을 생략했을 때 사용할 기본 자료형을 정한다.
    \item 템플릿 타입 결정과 생성자 인수의 암시적 형 변환은 별개의 단계이다.
    \item STL의 컨테이너는 클래스 템플릿이므로 Day 9의 내용이 STL 학습의 기반이 된다.
\end{itemize}
\end{importantbox}

\section{개념별 실습 완성 코드}
\subsection{함수 템플릿}
\begin{practicebox}
두 값을 곱하여 반환하는 함수 템플릿 \texttt{multiply()}를 작성한다. \texttt{int}와 \texttt{double} 자료형으로 각각 호출한다.
\end{practicebox}

\begin{lstlisting}
#include <iostream>

using namespace std;

template <typename T>
T multiply(T a, T b)
{
    return a * b;
}

int main()
{
    cout << multiply(3, 4) << endl;
    cout << multiply(2.5, 4.2) << endl;

    return 0;
}
\end{lstlisting}

\subsection{여러 개의 템플릿 타입}
\begin{practicebox}
서로 다른 두 자료형의 값을 한 줄에 출력하는 함수 템플릿 \texttt{printPair()}를 작성한다.
\end{practicebox}

\begin{lstlisting}
#include <iostream>

using namespace std;

template <typename T, typename U>
void printPair(T first, U second)
{
    cout << first << " " << second << endl;
}

int main()
{
    printPair("Kim", 90);
    printPair(3.5, 'A');

    return 0;
}
\end{lstlisting}

\subsection{일반 함수와 템플릿 함수}
\begin{practicebox}
\texttt{double} 일반 함수와 템플릿 함수를 같은 이름으로 작성하고, \texttt{double}과 \texttt{int} 인수에서 어떤 함수가 선택되는지 확인한다.
\end{practicebox}

\begin{lstlisting}
#include <iostream>

using namespace std;

void print(double x)
{
    cout << "Double Function" << endl;
}

template <typename T>
void print(T x)
{
    cout << "Template Function" << endl;
}

int main()
{
    print(3.14);
    print(100);

    return 0;
}
\end{lstlisting}

\subsection{클래스 템플릿}
\begin{practicebox}
하나의 값을 저장하고 출력하는 클래스 템플릿 \texttt{Data}를 작성한다. \texttt{string}과 \texttt{double} 자료형으로 객체를 생성한다.
\end{practicebox}

\begin{lstlisting}
#include <iostream>
#include <string>

using namespace std;

template <typename T>
class Data
{
public:
    T value;

    Data(T value)
    {
        this->value = value;
    }

    void show()
    {
        cout << value << endl;
    }
};

int main()
{
    Data<string> data1("Physics");
    Data<double> data2(4.24);

    data1.show();
    data2.show();

    return 0;
}
\end{lstlisting}

\subsection{기본 템플릿 인자}
\begin{practicebox}
기본 타입이 \texttt{double}인 클래스 템플릿 \texttt{Value}를 작성한다. 기본 타입과 직접 지정한 \texttt{char} 타입으로 객체를 생성한다.
\end{practicebox}

\begin{lstlisting}
#include <iostream>

using namespace std;

template <typename T = double>
class Value
{
public:
    T value;

    Value(T value)
    {
        this->value = value;
    }

    void show()
    {
        cout << value << endl;
    }
};

int main()
{
    Value<> v1(10);
    Value<char> v2('Z');

    v1.show();
    v2.show();

    return 0;
}
\end{lstlisting}

\section{실습 파일 구성}
\begin{practicebox}
Day 9 파일은 다음과 같이 관리한다.
\begin{itemize}[leftmargin=1.5em]
    \item \texttt{01\_function\_template.cpp}
    \item \texttt{02\_multiple\_types.cpp}
    \item \texttt{03\_template\_overloading.cpp}
    \item \texttt{04\_class\_template.cpp}
    \item \texttt{05\_template\_default.cpp}
    \item \texttt{practice/p01\_function\_template.cpp}
    \item \texttt{practice/p02\_multiple\_types.cpp}
    \item \texttt{practice/p03\_template\_overloading.cpp}
    \item \texttt{practice/p04\_class\_template.cpp}
    \item \texttt{practice/p05\_template\_default.cpp}
\end{itemize}
\end{practicebox}

\end{document}
